\SourceLeaf{286}%
können, ist zwar eine genäherte Kenntniss der Grösse der Linien, der orientirten Azimuthe, und der Breiten der Endpunkte erforderlich, was nur durch eine vorläufige Berechnung der Dreiecke zu erhalten ist: allein dieser Umstand ist durchaus unerheblich, da eine vorläufige schon die Ausführung der Messungen Schritt für Schritt begleitende Berechnung ohnehin in vielen Beziehungen räthlich, und zur Centrirung der excentrisch gemessenen Winkel, so wie zur Bestimmung des sphärischen oder sphäroidischen Excesses der Winkelsumme jedes Dreiecks sogar nothwendig ist; ja für den ersten Zweck wird, bei der Geringfügigkeit jener Reductionen, schon eine ganz rohe Annäherung immer zureichen, während das scharfe Centriren zuweilen, bei etwas beträchtlicher Excentricität der Standpunkte eine viel weiter getriebene Annäherung erfordern kann. Ich habe die Vorschriften deshalb entwickelt, damit man, \textit{wenn} man jene Reductionen berücksichtigen will, alles zu ihrer schärfsten Berechnung nöthige bereit finde, oder wenn man sie \textit{nicht} berücksichtigen will, leicht und bestimmt übersehen könne, wie wenig man dadurch aufopfert. Bei dem ganzen Hannoverschen Dreieckssystem sind die Reductionen durchgehends so äusserst gering, dass ihre Berücksichtigung als gänzlich überflüssig erscheint, und in der ganzen Ausdehnung der Zone von zwölf Breitegraden, für welche ich den Hülfsapparat beifüge, bleiben sie noch unterhalb derjenigen Bogensecundentheile, auf welche man sich bei den meisten Messungen in der Rechnung zu beschränken pflegt. Um diess recht evident hervortreten zu lassen, füge ich hier noch die numerische Rechnung für ein Paar Beispiele bei.

In dem Hannoverschen Dreieckssystem kommen die grössten Reductionen vor bei den Richtungen der Seiten des Dreiecks Brocken-Hohehagen-Inselsberg, welches Dreieck zugleich das grösste und das von dem Normal-Parallelkreise am entferntesten liegende ist; bei allen übrigen Dreiecksseiten überschreiten die Reductionen nirgends zwei Tausendtheile der Secunde, und die meisten erreichen nicht einmal den Werth \(0''001\).

Es ist für diese Punkte
\begin{center}
\begin{tabular}{l|c|c|r}
&\multicolumn{2}{c|}{Breite}&\\
&auf dem Ellipsoid&auf der Kugel&\(k\)\\\hline
Brocken&\(51^0\quad48'\quad2''\)&\(51^0\quad46'\quad3''\)&\(0''164\)\\
Hohehagen&\(51\quad28\quad31\)&\(51\quad26\quad35\)&\(0{,}303\)\\
Inselsberg&\(50\quad51\quad9\)&\(50\quad49\quad16\)&\(0{,}687\)
\end{tabular}
\end{center}
